%! AP Statistics — Unit 4, Lesson 4.4 — Activity KEY
\documentclass[10pt]{article}
\usepackage{apstats-article}
\usepackage{apstats-key}

\begin{document}

\pageheader{Unit 4, Lesson 4.4}{Group Activity KEY: Are These Events Mutually Exclusive?}
\namepartnerperiod

\vspace{0.1in}

\begin{scenariobox}[Context]{sky}
\small
A large high school surveyed 200 students about their after-school activities and
class standing. The two-way table below shows the results.

\vspace{6pt}
\renewcommand{\arraystretch}{1.4}
\begin{tabularx}{\linewidth}{@{} l *{4}{X} r @{}}
 & \textbf{Freshman} & \textbf{Sophomore} & \textbf{Junior} & \textbf{Senior} & \textbf{Total} \\
\hline
\textbf{Sports}    & 18 & 22 & 20 & 15 & 75  \\
\textbf{Arts/Music}& 12 & 10 & 14 & 9  & 45  \\
\textbf{Neither}   & 20 & 18 & 22 & 20 & 80  \\
\hline
\textbf{Total}     & 50 & 50 & 56 & 44 & 200 \\
\end{tabularx}
\end{scenariobox}

\vspace{0.12in}

\textbf{Tier R}
\begin{practicebox}
\small
\begin{enumerate}[leftmargin=1.5em, itemsep=6pt]
  \item \ans{No.} \quad \ans{ME} --- A student cannot be both a freshman and a senior simultaneously.
  \item \ans{Based on the table structure, each student is in exactly one row, so yes.} \quad \ans{ME}
  \item \ans{Yes --- 20 juniors play sports.} \quad \ans{Not ME}
  \item \ans{ME} --- The table separates Arts/Music and Neither into distinct rows; a student is in one or the other.
\end{enumerate}
\end{practicebox}

\vspace{0.08in}

\textbf{Tier A}
\begin{practicebox}
\small
\begin{enumerate}[leftmargin=1.5em, itemsep=8pt]
  \item
    \begin{enumerate}[label=(\alph*), itemsep=4pt]
      \item \ans{Not mutually exclusive. A student can be both a sophomore and in Sports --- there are 22 such students in the table.}
      \item \ans{$P(A \cap B) = 22/200 = 0.11$}
    \end{enumerate}

  \item
    \begin{enumerate}[label=(\alph*), itemsep=4pt]
      \item \ans{Mutually exclusive. A student cannot be both a freshman and a senior at the same time; a student is in exactly one grade level.}
      \item \ans{$P(A \cap B) = 0$}
    \end{enumerate}

  \item
    \begin{enumerate}[label=(\alph*), itemsep=4pt]
      \item \ans{Mutually exclusive (given the table design, each student is in exactly one activity row). No student appears in both the Sports row and the Arts/Music row.}
      \item \ans{$P(A \cap B) = 0$; $P(A) = 75/200 = 0.375$}
    \end{enumerate}
\end{enumerate}
\end{practicebox}

\vspace{0.08in}

\textbf{Tier E}
\begin{practicebox}
\small
\begin{enumerate}[leftmargin=1.5em, itemsep=8pt]
  \item \ans{Sample ME pair: Freshman and Senior. A student cannot be in both grade levels simultaneously. Sample not-ME pair: Junior and Sports. A junior can play sports (20 students in the table do), so the events share outcomes.}

  \item \ans{The reasoning is incorrect. Mutually exclusive is determined by whether events share any outcomes, not by the size of their probabilities. For example, "Freshman" ($P = 50/200 = 0.25$) and "Senior" ($P = 44/200 = 0.22$) are mutually exclusive even though neither probability is "small."}

  \item \ans{Counterexample: Flip a fair coin. $A = \{\text{heads}\}$, $B = \{\text{tails}\}$. Both have probability 0.5 --- not small at all --- and they are mutually exclusive because a coin cannot show both heads and tails on one flip.}
\end{enumerate}
\end{practicebox}

\end{document}
